\documentclass{beamer}

\usepackage[utf8]{inputenc}
\usepackage[T5]{fontenc}
\usepackage[vietnamese]{babel}

\usepackage{amsmath,amssymb}
\usepackage{tikz}
\usetikzlibrary{positioning,arrows.meta,calc,fit,shapes.geometric}
\usepackage{listings}
\usepackage{array}
\usepackage{colortbl}
\usepackage{booktabs}
\usepackage{xcolor}
\usepackage{tcolorbox}

\usetheme{Madrid}
\usecolortheme{dolphin}

% --- màu ---
\definecolor{myblue}{RGB}{30,144,255}
\definecolor{mygreen}{RGB}{0,180,0}
\definecolor{myred}{RGB}{220,20,60}
\definecolor{mypurple}{RGB}{138,43,226}
\definecolor{myorange}{RGB}{255,140,0}
\definecolor{mycyan}{RGB}{0,206,209}
\definecolor{mygray}{RGB}{90,90,90}

\setbeamercolor{structure}{fg=myblue}
\setbeamercolor{frametitle}{fg=white,bg=myblue}
\setbeamerfont{frametitle}{series=\bfseries}

\tcbset{
colback=blue!5,
colframe=myblue,
arc=3mm,
boxrule=0.8pt
}

% macros
\newcommand{\hlblue}[1]{\textcolor{myblue}{#1}}
\newcommand{\hlgreen}[1]{\textcolor{mygreen}{#1}}
\newcommand{\hlred}[1]{\textcolor{myred}{#1}}
\newcommand{\hlpurple}[1]{\textcolor{mypurple}{#1}}
\newcommand{\hlcyan}[1]{\textcolor{mycyan}{#1}}
\newcommand{\hlwhite}[1]{\textcolor{white}{#1}}

\title{\hlred{Bài tập DFS}}
\author{Group 3}
\date{\hlcyan{03/2026}}

\begin{document}

\frame{\titlepage}

% ===================== TỔNG QUAN =====================

\begin{frame}{\hlwhite{Tổng quan bài tập}}
\begin{tcolorbox}[title=\textbf{\hlwhite{Chủ đề}}]
Bài toán sử dụng \hlblue{DFS} để duyệt đồ thị và xác định thứ tự thăm các đỉnh.
\end{tcolorbox}

\begin{itemize}
\item Không dùng thuật toán nâng cao
\item Chỉ cần DFS cơ bản
\item Rèn luyện kỹ năng duyệt đồ thị
\end{itemize}
\end{frame}

% ===================== MINH HỌA =====================

\begin{frame}{\hlwhite{Minh họa}}
\begin{columns}
\column{0.6\textwidth}
\begin{tikzpicture}[scale=1]
\tikzstyle{v}=[circle,draw=myblue,thick,minimum size=8mm]

\node[v] (1) at (0,2) {1};
\node[v] (2) at (2,3) {2};
\node[v] (3) at (4,2) {3};
\node[v] (4) at (1,0.5) {4};
\node[v] (5) at (3,0.5) {5};

\draw[thick] (1)--(2);
\draw[thick] (2)--(3);
\draw[thick] (1)--(4);
\draw[thick] (2)--(5);
\draw[thick] (4)--(5);
\end{tikzpicture}

\column{0.4\textwidth}
\begin{tcolorbox}
\hlgreen{DFS sẽ đi sâu theo một nhánh trước}
\end{tcolorbox}
\end{columns}
\end{frame}

% ===================== ĐỀ BÀI =====================

\begin{frame}{\hlwhite{Đề bài}}
\begin{tcolorbox}[title=\textbf{\hlwhite{DFS Traversal}}]
Cho đồ thị vô hướng gồm $n$ đỉnh và $m$ cạnh.

Hãy thực hiện duyệt \hlblue{DFS} bắt đầu từ đỉnh $1$ và in ra thứ tự các đỉnh được thăm.
\end{tcolorbox}

\begin{itemize}
\item Các đỉnh được đánh số từ $1$ đến $n$
\item Nếu có nhiều lựa chọn, ưu tiên đỉnh có số nhỏ hơn
\end{itemize}
\end{frame}

% ===================== INPUT =====================

\begin{frame}{\hlwhite{Dữ liệu vào}}
\begin{tcolorbox}[title=\textbf{\hlwhite{Input}}]
Dòng đầu: $n, m$

$m$ dòng tiếp theo: mỗi dòng gồm $u, v$
\end{tcolorbox}

\begin{itemize}
\item $1 \le n \le 1000$
\item $0 \le m \le \dfrac{n(n-1)}{2}$
\end{itemize}
\end{frame}

% ===================== OUTPUT =====================

\begin{frame}{\hlwhite{Dữ liệu ra}}
\begin{tcolorbox}[title=\textbf{\hlwhite{Output}}]
In ra thứ tự DFS bắt đầu từ đỉnh $1$
\end{tcolorbox}

\begin{itemize}
\item Các đỉnh cách nhau bởi dấu cách
\end{itemize}
\end{frame}

% ===================== SAMPLE =====================

\begin{frame}[fragile]{\hlwhite{Ví dụ}}
\begin{columns}

\column{0.48\textwidth}
\begin{tcolorbox}[title=\textbf{Input}]
\begin{lstlisting}
5 5
1 2
2 3
1 4
2 5
4 5
\end{lstlisting}
\end{tcolorbox}

\column{0.48\textwidth}
\begin{tcolorbox}[title=\textbf{Output}]
\begin{lstlisting}
1 2 3 5 4
\end{lstlisting}
\end{tcolorbox}

\end{columns}
\end{frame}

% ===================== GIẢI THÍCH =====================

\begin{frame}{\hlwhite{Giải thích}}
\begin{itemize}
\item Bắt đầu từ 1 → đi 2
\item từ 2 → đi 3
\item quay lại → đi 5
\item từ 5 → đi 4
\end{itemize}

\begin{tcolorbox}
\centering
\hlpurple{DFS luôn đi sâu trước rồi mới quay lui}
\end{tcolorbox}
\end{frame}

% ===================== GỢI Ý =====================

\begin{frame}{\hlwhite{Gợi ý}}
\begin{itemize}
\item Dùng \hlblue{vector<int>} để lưu danh sách kề
\item Sắp xếp danh sách kề
\item Dùng DFS đệ quy
\item Mảng visited để tránh lặp
\end{itemize}

\begin{tcolorbox}
\centering
\hlred{Nếu không sort, thứ tự DFS sẽ sai}
\end{tcolorbox}
\end{frame}

% ===================== SUBTASK =====================

\begin{frame}{\hlwhite{Subtask}}
\begin{tabular}{|c|c|c|}
\hline
Subtask & Giới hạn & Điểm \\
\hline
1 & $n \le 20$ & 30\% \\
\hline
2 & $n \le 200$ & 30\% \\
\hline
3 & $n \le 1000$ & 40\% \\
\hline
\end{tabular}
\end{frame}

% ===================== KẾT =====================

\begin{frame}{\hlwhite{Lưu ý}}
\begin{itemize}
\item DFS phụ thuộc thứ tự duyệt
\item luôn cần visited
\item nên sort adjacency list
\end{itemize}

\begin{tcolorbox}
\centering
\hlblue{Đây là bài nền tảng cực quan trọng}
\end{tcolorbox}
\end{frame}

\end{document}